\chapter{Uso de comandos vim en LazyVim}
Vim es un editor de texto bastante popular que es muy usado para programar por la cantidad de complementos que se desarrollaron a su alrededor y hacen que esta pueda convertirse en un IDE profesional y no tenga nada que envidiar a los mas usados en el mundo tanto así que VisualStudio Code implemento plugins para convertir su IDE en uno similar a los muchos que tenemos en VIM.
\\
Vim soluciona algo que a muchos programadores les encanta y es NO USAR EL MOUSE para ganar rapidez, puede parecer algo raro pero una vez puedes hacer todo con las manos en el teclado no hay marcha atrás y aunque usar vim no se traduce en escribir mejor código directamente, es muy satisfactorio y rápido saber usarlo.
\\
Esta guía es un complemento a la instalación de LazyVim que realice. Repasaremos como usar los comandos mas usados para poder navegar y escribir código.
\section{Comandos para guardar, salir}
Cuando pulsamos : en el modo normal se nos abre una miniventana para comandos cmdline
\begin{itemize}
  \item w, guardar cambios.
  \item q, salir de vim.
  \item q!, salir de vim sin guardar cambios.
  \item wq, guardar y salir de vim.
\end{itemize}
\section{Modos VIM}
Existen 3 modos en VIM
\begin{enumerate}
  \item \textbf{normal}, pulsamos la tecla ESC, la tecla escape nos lleva a este modo lugar donde podremos navegar.
  \item \textbf{Insert}, pulsamos la tecla i, modo donde podremos escribir o insertar texto, primero tenemos que estar en el modo normal para entrar en este modo(Mas adelante veremos otros atajos para entrar en modo insert)
  \item \textbf{Visual}, pulsamos la tecla v, modo donde podremos seleccionar texto, primero tenemos que estar en el modo normal para entrar en el modo visual
\end{enumerate}
\section{Navegación}
\subsection{Navegación básica hjkl}
Para navegar en Vim usamos las teclas hjkl.
\begin{itemize}
  \item h, nos movemos a la izquierda.
  \item j, nos movemos abajo.
  \item k, nos movemos arriba.
  \item l, nos movemos a la derecha.
\end{itemize}
\subsection{Navegación PageUp PageDown}
Al igual que las teclas PageUP y PageDown tradicionales LazyVim tenemos los comandos:
\begin{itemize}
  \item Ctrl+u, PageUp
  \item Ctrl+d, PageDown
\end{itemize}
\subsection{Navegación numérica}
Usando la navegación básica hjkl podemos especificar cuantos lineas o caracteres nos queremos mover.
\begin{itemize}
  \item 12h, se moverá 12 caracteres a la izquierda.
  \item 110j, se moverá 110 lineas a la abajo.
  \item 5k, se moverá 5 lineas a la arriba.
  \item 4l, se moverá 4 caracteres a la derecha.
\end{itemize}
\subsection{Navegación inicio y final}
Para ir al inicio o final
\begin{itemize}
  \item {[ [}, nos vamos a la primera linea del archivo.
  \item ] ], nos vamos a la ultima linea del archivo.
\end{itemize}
\subsection{Navegación palabras}
En lugar de usar llll varias l para ir a la derecha o hhhh para ir a la izquierda, podemos avanzar o retroceder de palabra con:
\begin{itemize}
  \item w, avanzamos a la primera letra de la siguiente palabra.
  \item e, avanzamos a la ultima letra de la siguiente palabra.
  \item b, retrocedemos a la primera letra de la anterior palabra.
\end{itemize}
\section{Insertar}
Al modo insert se puede acceder de las siguientes maneras.
\begin{itemize}
  \item i, entra en modo insert en el lugar donde estemos.
  \item o, modo insert abajo de donde estamos creando una nueva linea.
  \item O, modo insert arriba de donde estamos creando una nueva linea.
  \item shift+a, modo insert al final de la linea.
  \item shift+i, modo insert al inicio de la linea.
\end{itemize}
\section{Visual, selecciona texto y párrafos}
Para seleccionar cosas para poder copiarlas y luego pegarlas primero tenemos que seleccionarlas con el modo visual muy similar cuando arrastramos el mouse para seleccionar varias lineas de texto aunque aquí no usamos el mouse.
\subsection{Visual, selección con hjkl}
\begin{itemize}
  \item vh, primero entramos en modo visual con v, luego navegamos con hhhh para poder seleccionar navegando hacia la izquierda.
  \item vjjjjj, primero entramos en modo visual con v, luego navegamos con jjjj para poder seleccionar navegando hacia abajo
  \item vkkkkk, primero entramos en modo visual con v, luego navegamos con kkkk para poder seleccionar navegando hacia arriba 
  \item vjjjjj, primero entramos en modo visual con v, luego navegamos con jjjj para poder seleccionar navegando hacia derecha.
\end{itemize}
\subsection{Visual, selección de palabras}
\begin{itemize}
  \item ve, Entramos en modo visual al inicio de la palabra luego con e seleccionamos hacia el final de la palabra
  \item veee, Entramos en modo visual al inicio de la palabra luego con eee vamos seleccionando hacia el final de la fila
  \item vb, Entramos en modo visual al final de la palabra luego con b seleccionamos hacia el inicio de la palabra
  \item vbbb, Entramos en modo visual al final de la palabra luego con bbb vamos seleccionando palabras hacia el inicio de la fila
\end{itemize}
\section{Ventanas}
En LazyVim podemos crear varias ventanas donde podemos mostrar los archivos que necesitemos y navegar entre ellas.
\subsection{Crear ventanas}
\begin{itemize}
  \item sv, creara una ventana vertical a la derecha,
  \item ss, creara una ventana horizontal abajo,
\end{itemize}
\subsection{Navegar entre ventanas}
\begin{itemize}
  \item sh, saltamos a la ventana izquierda 
  \item sj, saltamos a la ventana de abajo
  \item sk, saltamos a la ventana arriba
  \item sl, saltamos a la ventana derecha
\end{itemize}
\section{Consola}
Podemos abrir y cerrar la consola con \textbf{crtl+/}
\section{Buscador}
Podemos buscar archivos y texto.
\subsection{Buscador de archivos}
\begin{itemize}
  \item Esp Esp, para abrir el buscador de archivos en mi carpeta de trabajo.
  \item Esp f f, en el root dir.
  \item Esp f F, cwd lugar desde donde ejecutamos LazyVim.
\end{itemize}
\subsection{Buscador de texto}
\begin{itemize}
	\item Esp s g, busca en el root dir.
	\item Esp s G, busca en el lugar desde donde ejecutamos LazyVim.
\end{itemize}
\subsection{Buscador grep en el archivo}
\begin{itemize}
  \item Presionamos / y se abre un buscador.
  \item Para navegar entre las palabras coincidencias encontradas apretamos n para avanzar y N para retroceder.
\end{itemize}
\section{Copiar palabras, textos, párrafos}
Esto se combina con otros comando de selección de la parte visual
\begin{itemize}
  \item y, para copiar lo seleccionado
\end{itemize}
\section{Eliminar}
\begin{itemize}
  \item dd, para cortar toda una linea, se guardara en el porta papel para ser pegado.
  \item D, para cortar desde el puntero en adelante o a la derecha.
\end{itemize}

\section{Crear un nuevo archivo}
Creamos un nuevo archivo en el lugar donde se encontraba antes de ejecutar el comando nvim.
\begin{itemize}
  \item te, se abrirá la consola de comandos con tabedit nombre\_archivo.txt
  \item tab, para cambiar de ventana.
\end{itemize}
\section{Comentar código}
gcc, comentamos código.





